\chapter{Trabajo futuro}
La validación prospectiva comprueba tres mecanismos de agrupación y dos intervenciones de retorno bajo condiciones acotadas. Queda estudiar otros grafos y restricciones de reutilización para determinar cuándo la cota es alcanzable. El contraejemplo analítico sugiere distinguir la referencia por máquina del mínimo factible bajo aislamiento de jobs. No basta extrapolar desde las configuraciones con DME igual a uno.

Otro paso es comprobar tamaños y cargas mayores, manteniendo explícita la frontera de archivos completos y separando el cambio de ubicación del de contratos. Una evaluación temporal requeriría controlar carga y cachés y distinguir el efecto de los jobs del volumen entregado. También queda aislar el costo de generar registros nativos, pendiente de OE4. Los tiempos posteriores de exportación y cálculo no lo miden.
